% ============================================================
%  Results Section — Robustification and Adaptive Pricing
%  Paste into your paper's \section{Computational Study}
% ============================================================

% ---- Preamble additions needed in your main .tex file -----
% \usepackage{booktabs}
% \usepackage{multirow}
% \usepackage{siunitx}   % for S column type
% \usepackage{amsmath}
% \usepackage{xcolor}
% \usepackage{caption}
% -----------------------------------------------------------

% ============================================================
\subsection{Value of Robustification}
\label{subsec:value_robustification}
% ============================================================

\subsubsection*{Experiment Design}

To quantify the benefit of the robust facility location plan~$\mathbf{x}^{\mathrm{rob}}$
relative to the nominal plan~$\mathbf{x}^{\mathrm{nom}}$, we solve both problems across
a full factorial grid of instance parameters: willingness-to-pay scale
$v \in \{0.50,\, 0.75,\, 1.00,\, 1.25\}$, congestion cost
$w \in \{10,\, 100,\, 1{,}000\}$, and uncertainty budget
$\Gamma \in \{1,\, 2,\, 3,\, 4\}$, yielding 48 combinations.
The nominal problem is a mixed-integer second-order cone programme (MISOCP)
solved to a 1\% MIP gap; the robust problem is solved via the
column-and-constraint generation (C\&CG) algorithm with the same tolerance.

For each combination, both solutions are evaluated under~$N = 30$
out-of-sample disruption scenarios sampled uniformly from the uncertainty
set~$\Xi(\Gamma)$, as well as under the individual worst-case disruption
found by the separation oracle (eq.~17 in the paper).
Throughout this subsection, \emph{Scenario~A} refers to the full-recourse
(adaptive pricing) evaluation, which is the natural complement to the
robust formulation.  The \emph{robustification value}
$\Delta^{\mathrm{rob}}_k = \pi^{\mathrm{rob}}_k - \pi^{\mathrm{nom}}_k$
is computed scenario-by-scenario over the same sample, making the
comparison free of sampling noise.

\subsubsection*{Price of Robustness}

Table~\ref{tab:price_of_robustness} reports the guaranteed worst-case
profit of the robust solution ($\pi^{\mathrm{rob}}_{\mathrm{LB}}$, the
C\&CG lower bound) alongside the nominal optimal profit
($\pi^{\mathrm{nom}}$) and the resulting \emph{price of robustness}
$\text{PoR} = (\pi^{\mathrm{nom}} - \pi^{\mathrm{rob}}_{\mathrm{LB}})\,/\,\pi^{\mathrm{nom}}$,
at $w = 10$ across all four levels of $v$ and $\Gamma$.

\begin{table}[htbp]
\centering
\caption{Price of Robustness (\%) at $w = 10$.
         PoR$= (\pi^{\mathrm{nom}} - \pi^{\mathrm{rob}}_{\mathrm{LB}})/\pi^{\mathrm{nom}}$.
         ``—'' denotes the corner case where the robust solution opens
         no facility ($\pi^{\mathrm{rob}}_{\mathrm{LB}} \approx 0$).}
\label{tab:price_of_robustness}
\setlength{\tabcolsep}{6pt}
\begin{tabular}{lrrrr}
\toprule
& \multicolumn{4}{c}{Uncertainty budget $\Gamma$} \\
\cmidrule(lr){2-5}
$v$-scale & 1 & 2 & 3 & 4 \\
\midrule
0.50 & 45.2\% & 75.9\% & —   & —    \\
0.75 & 27.8\% & 46.8\% & 61.8\% & 75.9\% \\
1.00 & 20.3\% & 35.7\% & 46.1\% & 57.4\% \\
1.25 & 16.3\% & 29.0\% & 36.5\% & 46.0\% \\
\bottomrule
\end{tabular}
\end{table}

The price of robustness is monotonically increasing in $\Gamma$ and
decreasing in $v$.  The former is expected: a larger uncertainty budget
demands a more conservative facility plan, which sacrifices nominal
performance.  The latter reflects the compounding effect of higher
willingness-to-pay: in more profitable markets the robust solution can
hedge disruption risk without as large a relative sacrifice, because the
revenue upside from serving customers is substantial regardless of the
capacity level chosen.

A noteworthy corner case arises at $v = 0.50$ and $\Gamma \geq 3$:
the C\&CG algorithm yields a trivial robust solution that opens no
facility, achieving a guaranteed worst-case profit of zero.  This is
rational — when the market margin is thin and the disruption budget is
large enough to reduce capacity to near zero, any open facility would
generate a guaranteed loss under the adversary's best play.
The robust model correctly identifies this no-operation regime and
avoids it.

\subsubsection*{Worst-Case Robustification Value}

The worst-case robustification value
$\Delta^{\mathrm{rob}}_{\mathrm{WC}} =
\pi^{\mathrm{rob}}_{\mathrm{WC}} - \pi^{\mathrm{nom}}_{\mathrm{WC}}$
compares the minimax profits of the two solutions, where each is
evaluated under its own worst-case disruption.
Table~\ref{tab:wc_rob_val} presents this quantity at $w = 10$.

\begin{table}[htbp]
\centering
\caption{Worst-case robustification value $\Delta^{\mathrm{rob}}_{\mathrm{WC}}$
         (in thousands) at $w = 10$, full-recourse (Scenario~A).}
\label{tab:wc_rob_val}
\setlength{\tabcolsep}{6pt}
\begin{tabular}{lrrrr}
\toprule
& \multicolumn{4}{c}{Uncertainty budget $\Gamma$} \\
\cmidrule(lr){2-5}
$v$-scale & 1 & 2 & 3 & 4 \\
\midrule
0.50 & 3.8  & 18.3  & 35.5  & 68.8  \\
0.75 & 0.0  & 18.3  & 35.3  & 73.8  \\
1.00 & 3.9  & 32.7  & 77.5  & 125.9 \\
1.25 & 64.9 & 132.9 & 220.0 & 276.4 \\
\bottomrule
\end{tabular}
\end{table}

The worst-case robustification value is uniformly non-negative — by
construction, the robust solution cannot be worse than the nominal at
its own worst-case scenario.
At $\Gamma = 1$ the gains are modest (near zero for $v = 0.75$, since
a budget of one disruption unit barely strains a well-diversified
facility plan).  However, the gains escalate sharply with $\Gamma$: at
$\Gamma = 4$ and $v = 1.25$, the robust solution delivers
\$276\,400 more profit than the nominal solution under their respective
adversarial scenarios — a 71\% improvement relative to the nominal
worst-case profit of \$138\,100.

\subsubsection*{Average-Case Performance and the Robustness--Efficiency Trade-off}

A conventional criticism of robust optimisation is that protection against
worst-case outcomes comes at the cost of average-case performance.
Figure~\ref{fig:rob_avg} illustrates this trade-off through the
average robustification value $\bar{\Delta}^{\mathrm{rob}} =
\mathbb{E}_\varepsilon[\pi^{\mathrm{rob}} - \pi^{\mathrm{nom}}]$
estimated over 30 out-of-sample scenarios.

The results reveal a nuanced picture.  For low willingness-to-pay
($v = 0.50$) the nominal solution outperforms the robust on average
across all $\Gamma$ values.  This is consistent with the corner-case
finding: for $\Gamma \geq 3$ the robust solution opens nothing, so
average profit is zero, whereas the nominal plan still captures
positive revenue in mild disruption scenarios.

For moderate to high willingness-to-pay ($v \geq 1.00$), the
trade-off reverses at sufficiently large $\Gamma$.  At $v = 1.00$,
$w = 10$, the robust solution earns \$19\,089 \emph{more} on average
than the nominal at $\Gamma = 3$, and \$15\,199 more at $\Gamma = 4$.
At $v = 1.25$ the advantage is even larger: \$30\,106 at $\Gamma = 3$
and \$29\,889 at $\Gamma = 4$.
The intuition is that in high-margin markets the robust plan invests in
a broader facility footprint that happens to serve average-case
disruptions efficiently, not merely the worst case.

\subsubsection*{Tail-Risk Reduction}

Table~\ref{tab:cvar_rob} reports the CVaR at the 5\% level
($\text{CVaR}_{5\%}$, i.e.\ the mean of the worst 5\% of the
out-of-sample profit distribution) for both solutions at $w = 10$.

\begin{table}[htbp]
\centering
\caption{CVaR$_{5\%}$ of OOS profit (thousands) for nominal and
         robust solutions at $w = 10$.
         Relative improvement = (rob $-$ nom)\,/\,|nom|.}
\label{tab:cvar_rob}
\setlength{\tabcolsep}{5pt}
\begin{tabular}{llrrrr}
\toprule
& & \multicolumn{4}{c}{$\Gamma$} \\
\cmidrule(lr){3-6}
$v$ & Solution & 1 & 2 & 3 & 4 \\
\midrule
\multirow{3}{*}{0.75}
  & Nominal & 130.5 &  82.9 &  28.7 &   5.2 \\
  & Robust  & 130.5 &  97.3 &  76.3 &  47.6 \\
  & Improv. &  0.0\% & +17.4\% & +165.7\% & +822.3\% \\
\addlinespace
\multirow{3}{*}{1.00}
  & Nominal & 218.9 & 154.9 &  79.9 &  47.9 \\
  & Robust  & 222.8 & 193.0 & 153.6 & 126.2 \\
  & Improv. & +1.8\% & +24.6\% &  +92.3\% & +163.5\% \\
\addlinespace
\multirow{3}{*}{1.25}
  & Nominal & 254.3 & 155.1 & 107.0 &  $-$5.5 \\
  & Robust  & 319.2 & 271.7 & 246.8 & 211.0 \\
  & Improv. & +25.5\% & +75.2\% & +130.7\% & $>$1000\% \\
\bottomrule
\end{tabular}
\end{table}

Tail-risk reduction is the most consistent benefit of robustification.
CVaR$_{5\%}$ improves in every instance where the robust solution is
non-degenerate, and the improvement accelerates with $\Gamma$.
At $v = 0.75$, $\Gamma = 4$, the nominal solution's CVaR$_{5\%}$ falls
to only \$5\,158 — barely above break-even — while the robust solution
maintains \$47\,574, an 822\% improvement.
At $v = 1.25$, $\Gamma = 4$, the nominal CVaR$_{5\%}$ turns negative
(\$-5\,479), indicating expected losses in the worst 5\% of scenarios;
the robust solution keeps CVaR$_{5\%}$ at \$211\,000.

With respect to the probability of loss, the nominal solution exhibits
non-zero loss probability as early as $\Gamma = 3$ for $v = 0.50$
(2\%) and rises to 36\% at $v = 0.50$, $w = 1{,}000$, $\Gamma = 4$.
In contrast, the robust solution records zero loss probability across
virtually all non-degenerate instances.

\medskip
\noindent\textbf{Summary.}
The robust solution delivers a clear worst-case guarantee — quantified
by the worst-case robustification value and the CVaR$_{5\%}$ improvement
— at the cost of a reduced nominal profit (the price of robustness).
Crucially, this trade-off is not universal: in high-margin markets
($v \geq 1.00$) the robust solution dominates the nominal on
\emph{both} worst-case and average-case metrics at sufficiently large
$\Gamma$, suggesting that robustification is not merely a risk-aversion
device but can also be a sound efficiency strategy.


% ============================================================
\subsection{Value of Adaptive Pricing}
\label{subsec:value_adaptive_pricing}
% ============================================================

\subsubsection*{Experiment Design}

To isolate the contribution of pricing flexibility, we compare two
recourse policies applied to the same first-stage plan under each
disruption scenario:
\begin{itemize}
  \item \textbf{Scenario~A — Adaptive pricing:} After disruption
    $\boldsymbol{\varepsilon}$ is realised, the provider jointly
    re-optimises allocation~$\mathbf{y}$ and prices~$\mathbf{p}$
    by solving the inner SOCP (Corollary~1, eq.~15).
    Prices adjust via $p_i = v_i(1 - \sum_j y_{ij})$ to reflect the
    reduced capacity.
  \item \textbf{Scenario~B — Fixed-price recourse:} Prices are locked
    at their pre-disruption values $\hat{p}_i = v_i(1 - \text{fill}_i^*)$,
    materialising effective demand $\lambda_i = \Lambda_i \cdot \text{fill}_i^*$.
    Only allocation is re-optimised (Remark~2, eq.~8).
\end{itemize}

The \emph{Value of Adaptive Pricing}
$\mathrm{VAP} = \pi^{\mathrm{ScenA}} - \pi^{\mathrm{ScenB}}$
is computed for each of the 30 out-of-sample scenarios and for the
worst-case disruption, separately for the nominal plan
($\mathrm{VAP}^{\mathrm{nom}}$) and the robust plan
($\mathrm{VAP}^{\mathrm{rob}}$).
The same 48-combination parameter grid as Section~\ref{subsec:value_robustification}
is used, enabling direct comparison.

\subsubsection*{VAP is Strictly Positive Under All Conditions}

Across all 48 parameter combinations, the average VAP is strictly
positive for both the nominal and robust solutions, and for both the
worst-case and average-case evaluations.
This confirms the theoretical result: adaptive pricing is a weakly
dominant recourse policy — the provider can always choose to hold prices
fixed, so the additional optimisation freedom cannot reduce profit.

Table~\ref{tab:vap_wc} reports the worst-case VAP for the robust
solution at $w = 10$.

\begin{table}[htbp]
\centering
\caption{Worst-case VAP for the robust solution
         $\mathrm{VAP}^{\mathrm{rob}}_{\mathrm{WC}}$ (thousands)
         at $w = 10$.}
\label{tab:vap_wc}
\setlength{\tabcolsep}{6pt}
\begin{tabular}{lrrrr}
\toprule
& \multicolumn{4}{c}{$\Gamma$} \\
\cmidrule(lr){2-5}
$v$-scale & 1 & 2 & 3 & 4 \\
\midrule
0.50 &   0.0 &  13.1 &   0.0 &  0.0 \\
0.75 &  16.0 &  25.5 &  25.4 & 33.9 \\
1.00 &  15.7 &  16.9 &  18.6 & 18.8 \\
1.25 &  32.7 &  16.6 &  21.1 & 66.3 \\
\bottomrule
\end{tabular}
\end{table}

\subsubsection*{VAP Increases with Disruption Severity}

The average VAP exhibits a clear monotone relationship with the
uncertainty budget $\Gamma$.
Table~\ref{tab:vap_pct} reports the average VAP as a percentage of the
average Scenario~A profit (i.e.\ relative to the adaptive-pricing
benchmark) at $w = 10$.

\begin{table}[htbp]
\centering
\caption{Average VAP as a percentage of Scenario~A profit at $w = 10$.
         Entries show nominal / robust percentages.}
\label{tab:vap_pct}
\setlength{\tabcolsep}{5pt}
\begin{tabular}{lllll}
\toprule
& \multicolumn{4}{c}{$\Gamma$} \\
\cmidrule(lr){2-5}
$v$-scale & 1 & 2 & 3 & 4 \\
\midrule
0.75 & 3.6\,/\,3.6\% & 6.7\,/\,6.8\% & 11.1\,/\,15.1\% & 19.4\,/\,24.4\% \\
1.00 & 4.1\,/\,4.0\% & 7.3\,/\,8.1\% & 11.3\,/\,11.4\% & 18.4\,/\,16.2\% \\
1.25 & 4.0\,/\,3.7\% & 5.0\,/\,3.4\% & 10.9\,/\,8.4\%  & 15.2\,/\,17.7\% \\
\bottomrule
\end{tabular}
\end{table}

At $\Gamma = 1$ the VAP represents roughly 3--4\% of adaptive-pricing
profit for all $v$ values, confirming that a limited disruption budget
leaves little room for pricing to compensate for capacity loss.
As $\Gamma$ increases the share rises steadily: at $\Gamma = 4$ and
$v = 0.75$, VAP accounts for 19.4\% of profit under the nominal plan
and 24.4\% under the robust plan.
This pattern is intuitive — heavier disruptions reduce available
capacity more severely, creating a larger wedge between what adaptive
and fixed-price recourses can achieve.

The robust solution consistently exhibits a \emph{higher} relative VAP
than the nominal at large $\Gamma$ (e.g.\ 24.4\% vs.\ 19.4\% at
$v = 0.75$, $\Gamma = 4$).
This occurs because the robust facility configuration concentrates
capacity at fewer, more disruption-resilient locations, so the residual
demand after disruption is more sensitive to how prices are set:
adaptive pricing can effectively suppress or stimulate demand at those
sites, while a fixed price cannot.

\subsubsection*{Effect of Congestion Cost}

Higher congestion cost~$w$ dampens the absolute value of VAP.
At $v = 0.75$, $\Gamma = 2$, the average VAP for the nominal solution
falls from \$9\,604 at $w = 10$ to \$4\,858 at $w = 1{,}000$.
With heavy congestion, both recourse policies — adaptive and fixed —
face the same structural constraint: routing additional demand to any
surviving facility incurs rapidly increasing delays.
The pricing lever therefore has less room to operate, since the marginal
cost of additional allocation is high regardless of the price set.
The relative VAP (as a fraction of adaptive-pricing profit) tells a
similar story, although the decline is more gradual because nominal
profits are also lower under high~$w$.

\subsubsection*{Combined Gain: Robustification and Adaptive Pricing Together}

The \emph{combined gain} measures the full benefit of simultaneously
using the robust plan and the adaptive pricing recourse, relative to the
worst-case baseline of using the nominal plan with fixed prices:
\[
  \Delta^{\mathrm{combined}} = \pi^{\mathrm{rob}}_{\mathrm{ScenA}}
  - \pi^{\mathrm{nom}}_{\mathrm{ScenB}}.
\]
Table~\ref{tab:combined} reports the average combined gain at $w = 10$.

\begin{table}[htbp]
\centering
\caption{Average combined gain $\bar\Delta^{\mathrm{combined}}$
         (thousands) at $w = 10$: robust+adaptive vs.\
         nominal+fixed-price.}
\label{tab:combined}
\setlength{\tabcolsep}{6pt}
\begin{tabular}{lrrrr}
\toprule
& \multicolumn{4}{c}{$\Gamma$} \\
\cmidrule(lr){2-5}
$v$-scale & 1 & 2 & 3 & 4 \\
\midrule
0.50 &   0.1 &  $-$6.0 & $-$37.6 & $-$17.8 \\
0.75 &   5.8 &  $-$9.5 &  11.4  &   5.0  \\
1.00 &   9.7 &  23.0   &  41.8  &  44.3  \\
1.25 &  15.8 & $-$12.4 &  60.9  &  67.1  \\
\bottomrule
\end{tabular}
\end{table}

The combined gain is positive in all moderate-to-high margin settings
($v \geq 1.00$) at $\Gamma \geq 1$, and consistently grows with
$\Gamma$ for $v = 1.00$ and $v = 1.25$.
At $v = 1.25$, $\Gamma = 3$, the provider who adopts both strategies
earns on average \$60\,900 more than one who uses neither — a gain of
approximately 21\% of the nominal (no-disruption) profit.

Negative combined gains appear primarily at $v = 0.50$ for large
$\Gamma$, driven by the previously noted corner case: the robust
solution opens no facility, so its adaptive-pricing profit is zero,
while the nominal solution still earns positive average revenue in
mild-disruption scenarios.
For $v \geq 0.75$ and $\Gamma \geq 3$, the combined gain is positive
in most configurations, indicating that the two strategies are
\emph{complements} rather than substitutes: robustification builds a
resilient capacity structure while adaptive pricing extracts maximum
value from whatever capacity survives the disruption.

\medskip
\noindent\textbf{Summary.}
Adaptive pricing delivers a strictly positive value under every
disruption scenario and for every facility plan tested.
The VAP ranges from 2--4\% of profit at low disruption ($\Gamma = 1$)
to 15--25\% at high disruption ($\Gamma = 4$), confirming that pricing
flexibility is most critical precisely when capacity is most stressed.
The robust facility plan amplifies this benefit: its concentrated,
disruption-resilient configuration makes the residual system more
responsive to price signals, yielding a higher relative VAP than the
nominal plan at large $\Gamma$.
Taken together, robust facility location and adaptive pricing are
complementary operational levers whose joint deployment can increase
expected profit by up to 21\% over the naive nominal-plan,
fixed-price baseline in high-margin, high-disruption settings.
